%=====================================================================
%  KERANGKA-SLIDE.TEX — PERPUSTAKAAN KERANGKA SLIDE SIAP PAKAI
%=====================================================================
%  CARA PAKAI:
%   1. Cari kerangka yang diinginkan di bawah (semuanya BERKOMENTAR).
%   2. Salin bloknya ke file konten Anda (03-... s.d. 08-...), atau
%      buat file baru lalu daftarkan di main.tex dengan \input{...}.
%   3. Hilangkan tanda % dan ganti teks tempat (---) dengan konten.
%  File ini TIDAK menghasilkan slide apa pun (isi 100% komentar).
%=====================================================================

%#####################################################################
%  A. SLIDE STANDAR (satu kolom)
%#####################################################################
%\begin{frame}
%    \frametitle{--- Judul Slide ---}
%    \justifying                      % rata kiri-kanan (opsional)
%    \small                           % ukuran font: \tiny \scriptsize \small
%    --- Tulis isi slide di sini. ---
%    % Footnote kutipan: \firstcite{kunci} (pertama) / \secondcite{kunci}
%    % Footnote biasa : teks\footnote{catatan kaki}
%\end{frame}

%#####################################################################
%  B. SLIDE DUA KOLOM (paling sering dipakai)
%#####################################################################
%\begin{frame}
%    \frametitle{--- Judul Slide ---}
%    \vspace*{-0.5cm}                 % hemat ruang vertikal (opsional)
%    \begin{columns}[T]               % [T] = rata atas
%        \begin{column}{0.5\textwidth}
%            \small
%            \begin{itemize}\tiny\justifying
%                \item Poin pertama
%                \item Poin kedua
%            \end{itemize}
%        \end{column}
%        \begin{column}{0.5\textwidth}
%            \small
%            --- isi kolom kanan ---
%        \end{column}
%    \end{columns}
%\end{frame}

%#####################################################################
%  C. SLIDE DENGAN GAMBAR + CAPTION + RUJUKAN SILANG
%#####################################################################
%\begin{frame}
%    \frametitle{--- Judul Slide ---}
%    \begin{figure}
%        \centering
%        \includegraphics[width=0.8\textwidth]{gambar/---folder/---berkas.png}
%        \captionsetup{font=tiny, justification=justified}
%        \captionof{figure}{--- keterangan gambar ---}
%        \label{fig:---label-unik---}          % rujuk: Gambar~\ref{fig:---}
%    \end{figure}
%\end{frame}

%#####################################################################
%  D. SLIDE DENGAN TABEL (booktabs)
%#####################################################################
%\begin{frame}
%    \frametitle{--- Judul Slide ---}
%    \begin{table}
%        \centering
%        {\small
%        \begin{tabular}{lcc}         % l=teks, c=tengah, r=kanan
%            \toprule
%            \textbf{Kolom 1} & \textbf{Kolom 2} & \textbf{Kolom 3} \\
%            \midrule
%            Baris A & isi & isi \\
%            Baris B & isi & isi \\
%            \bottomrule
%        \end{tabular}}
%        \captionof{table}{--- keterangan tabel ---}
%        \label{tab:---label-unik---}          % rujuk: Tabel~\ref{tab:---}
%    \end{table}
%\end{frame}

%#####################################################################
%  E. SLIDE PERSAMAAN MATEMATIKA / KUANTUM
%#####################################################################
%\begin{frame}
%    \frametitle{--- Judul Slide ---}
%    Satu persamaan berlabel (rujuk: Persamaan~\ref{eq:---}):
%    \begin{equation}
%        S(\rho) = -\Tr(\rho\log\rho)
%        \label{eq:---label-unik---}
%    \end{equation}
%    Beberapa baris tanpa nomor:
%    \begin{align*}
%        \ket{\psi} &= \alpha\ket{0}+\beta\ket{1}\\
%        \rho       &= \ketbra{\psi}{\psi}
%    \end{align*}
%\end{frame}

%#####################################################################
%  F. SLIDE KODE PROGRAM  (WAJIB [fragile]!)
%#####################################################################
%\begin{frame}[fragile]
%    \frametitle{--- Judul Slide ---}
%    \begin{lstlisting}[language=Python]     % ganti sesuai bahasa
%    --- kode di sini, indentasi jangan pakai tab ---
%    \end{lstlisting}
%\end{frame}

%#####################################################################
%  G. SLIDE DENGAN BLOK BERTITIK-TITIK (3 jenis)
%#####################################################################
%\begin{frame}
%    \frametitle{--- Judul Slide ---}
%    \begin{block}{Blok biasa}        % emas
%        --- definisi / catatan ---
%    \end{block}
%    \begin{alertblock}{Blok peringatan}  % oranye
%        --- hal PENTING / kesimpulan ---
%    \end{alertblock}
%    \begin{exampleblock}{Blok contoh}    % biru
%        --- contoh / studi kasus ---
%    \end{exampleblock}
%\end{frame}

%#####################################################################
%  H. SLIDE TANPA NOMOR (lampiran/cadangan)
%#####################################################################
%\begin{frame}[noframenumbering]
%    \frametitle{--- Slide Cadangan ---}
%    --- isi ---
%\end{frame}

%#####################################################################
%  I. SUB-BAGIAN (muncul sebagai titik kedua di navigasi atas)
%#####################################################################
%\section{--- BAGIAN BESAR ---}
%\subsection{--- Sub bagian ---}
%\begin{frame}
%    \frametitle{--- Judul Slide ---}
%    --- isi ---
%\end{frame}
